\documentclass[12pt]{article}

\usepackage[a4paper, total={6in, 8in}]{geometry}
\usepackage{textcomp}

\begin{document}
\setlength{\parindent}{0pt}

\def \t {\textrightarrow}
\def \v {\vspace{1em}}
\def \bi {\begin{itemize}}
\def \ei {\end{itemize}}
\def \s[#1] {\section*{#1}}
\def \ss[#1] {\subsection*{#1}}
\def \sss[#1] {\subsubsection*{#1}}

\s[Pirandello - teatro]
22 esce l'Enrico IV \t fondamentale è la pazzia, e di usarla a proprio vantaggio: non egoistico, ma esternazione di una scappatoia dalla piattezza della vita quotidiana

\ss[Uomo dal fiore in bocca]
"Uomo dal fiore in bocca" \t è un monoatto, che presente:
\bi
  \item giudizio pubblico e quello privato
  \item la realta dentro di noi e quella esterna
  \item conflittualità estrema, dovuta al non avere riferimenti \t atmosfera è fuori da ogni contesto socio-culturale
\ei
Uomo non viene nominato, e a un certo punto scopre di avere un tumore in bocca \t la sua vita ora cambia, e si ribaltano le dinamiche tra forma e vita \\
La forma è quella che viene attribuita dalla società, mentre la vita è quella che fluisce e che porta fuori in molte manifestazioni \t es. lo studio medico al quale si rivolge, viene osservato con occhi + critici \\
Scopre le forme delle sedie, del tavolo estremamente pulito e rivede una realtà gia nota con occhi diversi \\
Questo accade anche nello spazio aperto \t cerca la sua immagine nelle vetrine, e ha visione diversa di se \t si coglie in quegli ultimi attimi di vita \\
Vivere appieno proprio quando la vita cessa di esistere \t avere i giorni contati porta a concentrarsi su quello che non si è ancora fatto \\
Ritrova anche rapporto con la moglie, che rivede in modo diverso \\
Un uomo che vive secondo la vita interiore, il suo modo per aggrapparsi alla vita è di aggrapparsi alla vita vera, che non è la forma che incarcera \\
Somiglianza intensa con Ungaretti "non sono stato mai tanto attaccato alla vita di fronte alle brutture della guerra" \\
Nella parabola del teatro, si replica quello che si trova nei romanzi \t relativismo gnoseologico, lotta contro convenzioni sociali, visione paradossale dell'esistere (trovato il senso della vita, ora le deve perdere) \\
Segni misteri e tragici, mischiati ad altri comici ed umoristici \\
Il riso del clown \t che ha sempre angoli della bocca all'in giù, ed è truccato \t da cui esce però sempre venatura di tristezza \\
Hanno in comune la sofferenza celata dietro altro \t la sofferenza celata è la vita vera, mentre la forma è il sorriso ostentato \\
Hanno in comune la sofferenza celata dietro altro \t la sofferenza celata è la vita vera, mentre la forma è il sorriso ostentato

\v

Meta letteratura è la letteratura che parla di se \t esiste anche il cinema nel cinema \\
Nella sua parabola di vita l'ultima opera che torna preoptente anche se non compiuta nel contesto drammaturigico \t è il "Gigante della montagna" \\
Qua c'è conflittualità tra i personaggi e l'autore, che è giudicante nei confronti dei personaggi (questo succede anche nel teatro) \\
I romanzi di pirandello sono atemporali, hanno valore gnomico \t vizi dell'umanità sono trasversali e attraversano ogni epoca, non c'è uomo infelice in una sola epoca \\
Levarsi da una scena storica significa mettere a nudo l'animo di chi si descrive \t scompare la visione provvidenziale, che in Verga in qualche modo era presente (padron Toni, malavoglia, rappresenta un barlume di speranza terrena) \\
Assstiamo qua alle tragedie degradate \t il vivere anche qunado abbiamo la maschera di un personaggio storico, è smarrire se stessi \t enrico iv è morto infatti, ma anche in quel caso il vivere assume le stesse caratteristiche di ??? \\
Questo assume termini antitetici, dualismo \t che sono finzione e realtà

\v

L'espressionse del grottesco \t con Bertold Brecht, contesta il potere \t il grottesco evoca l'assurdo \\
Con l'espressione assurdità oggi si intende tutto ciò che è fuori dalla realta, e si vede in un certo contesto però \\
E quindi ribaltamento \t tutto ciò che è ordinario viene ribaltato \t è diverso da straordinario, che è possibile che succede, mentre assurdo è fuori dai parametri del reale \\
Tutto questo si coglie in un clima che devasta la società di quei decenni \t era normale che si sentisse passione o commozione per uno spargimento di sangue? \\
Si parla di new age, con le pratiche del se interiore \t tanto attesa quanto le guerre ?? \\
Pirandello da anche report tecnologico dell'epoca \t con l'invenzione della luce elettrica, che è garanzia di un lavoro che può seguire un espasione temporale più ampia della gioranta \\
Città popolose, chiassose, inquinate \t i primi ventenni del movimento stra paese stra citta

\v

"Quaderni" e "Ciack si gira" \t la gestazione lunghissima \t rientrano in una società che si allinea al progresso \\
Attualità estrema \\
Quaderni \t pubblicata nel 1915 e 25 (con modifiche nel titolo, che da "Ciack si gira" diventa i "Quaderni di serafino abb?") \\
È l'emblema dello sperimentalismo dell'autore \t è un romanzo aperto, le vicende vanno lette dal lettore, interpretate e completate \t una frammentazione molto forte, ancora più aperta di il fu mattia \\
Serafino bubbio \t insegnante disoccupato, che lavora presso un agenzia cinematografica \t e vede immischiato nella macchina della società moderna, che avvilupa e riduce l' uomo a un marchingegno \t spersonalizza \\
Nella qua abitudinarietà una troup di lavoro doveva girare una scena di un film, in cui una tigre doveva rivoltarsi nei confronti della donna \t l'attrice è una figura importante, è un avventuriera \\
Si era gia distinta per essere stata responsabile della fine di vita di un allievo \t partecipa nella scena, e l'attore uomo doveva dibattere con lei \\
La tigre avrebbe improvvisamente aggredito la donna e lacerata \t lui assiste a questa scena e perde la parola, questa Spannung \\
Lui diventa muto, e questo mostra la violenza del progresso nei confronti dell'uomo \t è il mutismo che hanno vissuto l'Aiala, che ha vissuto il Mostrada, tutti i personaggi \t il mutismo davanti alla prepotenza della societa, che sovrasta l'identita e la distrugge \\
La spersonalizzazione dell'uomo è nella spersonalizzazione della macchina \t la covolesis che descrive il mondo proprio presentato, dall'universo di James Joyce \\
Il mondo di dublino, paralizzata e soffocata dalla nebbia \\
Il suo mutismo è il silenzio delle cose, che porta a sviluppare un ascolto \t è il mutismo dell'intellettuale di fronte al potere ridotto a una mera obbedienza, come l'uomo è sottoposto alla mera gestione della macchina = degrado dell'intellettuale moderno \\
In parallelo l'intelligenza generativa \t il silenzio è l'unico strumento di denuncia \t il mondo dello spettacolo è compreso, e descrive il fallimento dell'assertività teatrale \\
Cinematografia ha svilito il teatro e ha cessato di avere rapporto accudente e formativo nei confronti di un pubblico ? \t l'opera è una critica alla modernizzazione \\
L'umorismo si allontana dal comico e si avvicina al pessimismo \t un pessimsimo che è la sfiducia nel progresso, e in noi stessi \\
L'opera è un paradigma, ed è l'epifany

\v

Gramsci, i quaderni dal carcere \t sono manifestazione di una scena politica che lo confina perchè in opposizione ai trend politici del tempo \\
Essere intellettuale significa prestarsi all'insubordinazione \t vista con Tacito, machiavelli etc \\
Alienazione, e riprodurre un lavoro senza creatività diventa alienante che impedisce autenticita \\
La maschera è presente in ognuno di noi

\v

Leone de Castris \t "Saggio sul decadentismo" \t il decadentismo italiano (svevo pirandello d'annunzio), testimonianza dell'importanza di pirandello a questa sede \\
L'inetto \t lui trasla la sua vita nelle sue opere \t ha una dobbia vita, e i suoi personaggi sono maldestri \t c'è la voglia di uscire dall'inettitudine? \\
Adriano Melci prova che il vecchio mattia pascal cerca un altra vita, un altro amore \t non chiude per sempre \t i personaggi di Svevo invece no \\
I suoi personaggi sono inetti, si lasciano andare in questo piattume \t neanche questa coscienza
\end{document}
